\section{Conclusion}
\label{sec:conclusion}

This project studied volatility as a tradable asset class through variance swaps, model-implied VIX derivatives, volatility-surface analysis, calibration, and volatility risk premium backtesting. The main objective was to connect the theoretical structure of volatility-linked products with working numerical methods and empirical analysis.

The first part of the project reviewed the role of volatility derivatives in modern financial markets. Variance swaps, VIX futures, and VIX options provide direct exposure to volatility and variance rather than only to the direction of the underlying asset. This motivates a modelling framework in which variance is treated as a stochastic state variable. The Heston model was used as the baseline framework because its variance process is mean-reverting, non-negative in continuous time, and analytically tractable enough for both simulation and characteristic-function-based pricing.

The theoretical part derived the key formulas required for the numerical implementation. The fair variance swap strike was obtained from the expected integrated variance under the Heston/CIR variance process. The project also constructed a model-implied VIX proxy based on the conditional expectation of future variance. This distinction is important because the proxy used in the project is not the official Cboe VIX index. Instead, it is a Heston-implied quantity of the form
\[
\widetilde{\mathrm{VIX}}_T
=
100\sqrt{A_\Delta+B_\Delta v_T}.
\]
This representation allowed VIX-type payoffs to be expressed as functions of terminal variance.

The W3 component implemented the Monte Carlo engine. It simulated Heston stock and variance paths using a full-truncation Euler scheme, computed realised variance from squared log returns, and compared the simulated realised variance with the analytical Heston variance swap strike. This comparison served as the first validation step for the simulation engine. The same Monte Carlo framework was also used to price a model-implied VIX call by transforming terminal variance into VIX and averaging discounted payoffs.

The W4 component implemented the CIR-based COS pricing engine. The method used the characteristic function of terminal CIR variance and approximated the expected payoff through a Fourier-cosine expansion. This provided a deterministic semi-analytical benchmark for the Monte Carlo VIX call price. The COS framework was also used for density recovery, VIX call and put pricing, VIX futures term-structure calculation, and convergence checks with respect to the number of terms and truncation interval.

The numerical comparison between the Monte Carlo and COS engines showed close agreement under the same parameter set. The reported COS price was approximately \(2.008969\), while the Monte Carlo price was approximately \(1.984952\). The absolute difference was \(0.024017\), corresponding to a relative difference of about \(1.21\%\). This result supports the internal consistency of the two independently implemented pricing methods.

The W5 component extended the project to market data. It introduced VIX index data, SPX market data, VIX futures data, realised variance estimation, implied variance approximation, volatility risk premium computation, and crisis-period analysis. The empirical findings confirmed several common volatility-market patterns: VIX rises during stress, implied variance is often higher than future realised variance in calm regimes, and volatility risk premium behaviour changes during crisis periods.

The W6 component implemented joint SPX--VIX calibration. The results showed a structural limitation of the one-factor Heston model. SPX-only calibration fitted the SPX implied-volatility surface better, while VIX-only calibration fitted the VIX futures curve better. The joint calibration did not dominate both single-market calibrations. This confirms that one variance factor is often insufficient to fit both SPX and VIX markets simultaneously.

The W7 component analysed volatility surfaces. The comparison between VIX and SPX smiles showed that VIX options tend to exhibit positive skew, while SPX options tend to exhibit negative skew. This difference has both mathematical and economic explanations. Under the CIR model, terminal variance is positive and right-skewed, which supports the right-skewed VIX distribution. In SPX options, negative stock-variance correlation and demand for downside protection support the left-skewed volatility smile.

The W8 component evaluated volatility risk premium backtesting and risk metrics. The results showed that high-VRP regimes can produce attractive risk-adjusted performance, but the full-period results also reveal substantial tail risk through drawdowns, negative skewness, and high kurtosis. This confirms that volatility-selling strategies cannot be evaluated only by average return or Sharpe ratio. Tail risk, crisis exposure, and regime dependence are central to the economic interpretation.

The main contribution of the project is therefore a complete research pipeline from theory to implementation and empirical testing. The project includes literature review, mathematical derivations, Monte Carlo simulation, COS pricing, data collection, calibration, volatility-surface analysis, and VRP backtesting. The work also includes a modular codebase, notebooks, diagnostic outputs, and a test suite that validates the main numerical components.

Several limitations remain. First, the official Cboe VIX methodology is more complex than the model-implied VIX proxy used in the theoretical and pricing sections. Second, the one-factor Heston/CIR model cannot fully capture the joint behaviour of SPX and VIX markets. Third, some empirical components rely on available data pipelines, cached snapshots, or synthetic fallback surfaces when live data are unavailable. Fourth, the COS method depends on truncation intervals and numerical payoff integration. Finally, volatility risk premium backtests remain sensitive to regime definitions, transaction costs, and crisis-period assumptions.

Future work should focus on improving both the model and the empirical strategy design. On the modelling side, possible extensions include two-factor CIR models, forward-variance models, rough Heston dynamics, and jump-diffusion specifications. These models may provide more flexibility for fitting SPX smiles, VIX futures, and VIX option surfaces simultaneously. On the numerical side, closed-form COS payoff coefficients and adaptive truncation rules could improve efficiency and accuracy. On the empirical side, future work should incorporate live market data, transaction costs, robust hedging rules, and more detailed stress testing.

Overall, the completed project demonstrates that volatility can be studied as a tradable asset class through a combination of stochastic modelling, numerical pricing, calibration, surface analysis, and backtesting. The Monte Carlo and COS engines provide the numerical foundation, the calibration results reveal the limitations of a one-factor model, the volatility-surface analysis explains the difference between VIX and SPX options, and the VRP backtest connects the modelling framework to investment and risk-management implications.